% sec-10: what is not claimed, and what would falsify the case.  Figure 20.
\section{Risks and Limits}

\subsection{What Is Not Claimed}

This section is written against the paper rather than for it. Each row below
pairs an inference a reader might reasonably draw with what is actually true.

\begin{apptable}
\begin{tabularx}{\textwidth}{>{\raggedright\arraybackslash}p{5.6cm} >{\raggedright\arraybackslash}X}
\headrow \thead{Inference a reader might draw} & \thead{What is actually true}\\
\midrule
The simulations validate RASolute 302 & They preceded it. The proximity is a chronology observation, and three differences are material: 1000 simulated against 241 enrolled, a combination against a single agent, KRAS G12C against a primarily G12D and G12V population \\
Daraxonrasib is first-in-human & It is investigational and already in Phase 3 evaluation. The supportable claim concerns the integrated surgical and advisory workflow \\
A robotic configuration has been specified or cleared & None has. Manufacturer, model, cleared configuration, arm count, instruments and software version are named at the site agreement \\
Drug supply or a cross-reference exists & Neither exists. No agreement of any kind is in place with the agent's developer \\
UC San Diego has agreed to something & It has not. One application asks for a 45-minute feasibility meeting \\
A patient has been treated & None has. The trial does not exist until a site, an IRB, an active IND, monitoring, insurance, and a sponsor of record are all in place \\
The advisory system holds a regulatory clearance & It holds none, and the applications ask for the study that would generate the evidence a clearance would require \\
\end{tabularx}
\end{apptable}
\tabcap{Seven inferences a reader might reasonably draw, and what is actually
true in each case. Every row corresponds to a sentence carried in the back
matter of all ten applications.}

\subsection{What Would Falsify the Case}

A claim that cannot fail is not a claim. Two released fields can falsify the
programme's central assertion, and both are released whatever they show.

\begin{appfloat}[!tb]
\begin{appfig}[graphviz-type / record nodes]
% Five record columns, each a stack of five anchored cells 2.7cm wide and 0.35
% apart.  No edges: joins are figure 15's subject, and an edge here would import
% the wrong question.
\foreach \i/\x/\ttl/\fa/\fb/\fc/\q/\hl in {%
  1/0/{stop\_latency\_record}/{arm\_id : int}/{measured\_ms : float}/{bench\_or\_intraop}/{Is the interlock as fast as claimed?}/gvcell,
  2/2.9/{advisory\_decision}/{case\_id : uuid}/{accepted : bool}/{latency\_ms : float}/{How often is advice rejected?}/gvcell,
  3/5.8/{escalation\_cohort}/{dose\_level : int}/{dlt\_count : int}/{negative\_results}/{Were negatives reported?}/gvcells,
  4/8.7/{simulation\_credibility}/{test\_id : int}/{passed : bool}/{credibility\_score}/{Does 81.9 reproduce?}/gvcells,
  5/11.6/{provenance\_record}/{artifact\_id : uuid}/{model\_commit : sha}/{reviewer : enum}/{Who checked this, and when?}/gvcell}{%
  \node[gvcellh,text width=24mm,minimum height=5.4mm] (h\i) at (\x,0) {\ttl};
  \node[gvcell,text width=24mm,minimum height=5.0mm,anchor=north] (a\i) at (h\i.south) {\fa};
  \node[gvcell,text width=24mm,minimum height=5.0mm,anchor=north] (b\i) at (a\i.south) {\fb};
  \node[\hl,text width=24mm,minimum height=5.0mm,anchor=north] (c\i) at (b\i.south) {\fc};
  \node[gvcellg,text width=24mm,minimum height=7.6mm,anchor=north,font=\tiny\itshape] (d\i) at (c\i.south) {\q};}
\node[font=\tiny\itshape,text=protogray,anchor=north,text width=128mm,align=center]
  at (5.8,-3.3) {The two tinted fields are the ones that can falsify the
  programme. No edges are drawn: joins are the subject of the data-schema
  figure, and an edge here would import the wrong question.};
\end{appfig}
\vspace{-0.7cm}
\figcaption{Five verification artifacts and the reviewer question each field
answers. Every field is released with the cohort, and the two marked fields are
the ones that can falsify the programme's central claim.}
\end{appfloat}

The first falsifier is \texttt{credibility\_score}. The digital-twin work carries
a score of 81.9 under ASME V\&V 40 and ICH M15 alignment across 55 verification
notebook tests \cite{daraxsims,asmevv40,ichm15}. If an independent run of those
tests does not reproduce that score, the modelling tier of the evidence stack
loses its standing, and with it the arm-selection rationale.

The second is \texttt{negative\_results}. The applications commit to releasing
negative and inconclusive cohorts, and \S6's Gold Standard Science mapping makes
that a field rather than a promise \cite{eo14303goldstandard}. A released cohort
with that field empty falsifies the programme's claim to Gold Standard alignment
directly, whatever the clinical result was.

\subsection{Risks the Programme Cannot Retire Alone}

\begin{apptable}
\begin{tabularx}{\textwidth}{>{\raggedright\arraybackslash}p{3.8cm} >{\raggedright\arraybackslash}p{3.4cm} >{\raggedright\arraybackslash}X}
\headrow \thead{Risk} & \thead{Who can retire it} & \thead{What the applicant can do meanwhile}\\
\midrule
No qualified site & An academic cancer center & Keep the regulatory package with the performer, so a second site substitutes without a protocol change \\
No drug supply or cross-reference & The agent's developer & State plainly in every document that none exists, and design the request so it can be answered in one letter \\
No sponsor-investigator determination & Institutional contracting and regulatory offices & Ask for the determination before the IRB submission, not after \\
Robotic configuration not permissible & The site's surgical robotics programme & Leave the configuration unspecified until the site names it \\
Accrual below plan & The site's referral base & Extend the accrual window rather than widening eligibility \\
\end{tabularx}
\end{apptable}
\tabcap{Five risks the programme cannot retire on its own, who can retire each,
and the action available meanwhile. None of the five is retired by funding,
which is why the funding argument and the feasibility argument are separate.}

\subsection{The Limit of the Whole Exercise}

Ten applications on one piece of science demonstrate that the mechanism is a
variable an individual applicant can address. They do not demonstrate that the
science is right, that the operation is safe, or that the drug helps in the
resectable setting. Those questions are settled by the trial, and the trial does
not exist. What this paper reports is the state of the argument as of August 3,
2026, with every number traceable to a named file and every claim paired with
the limitation its authors stated.
